\chapter[The Fixed Normalized Generative Law]{The Fixed Normalized\\Generative Law}
\label{ch:generative}

This chapter fixes one normalized joint before inference begins. The latent
channels inherit the two representations of the common principal gauge group
from \Cref{ch:geometry}. Their local coordinates may be re-expressed
independently, so the covariance formulas carry a passive product of frame
changes without introducing a second principal bundle. The linear-Gaussian
realization is deferred to \Cref{ch:generative-gaussian}.

\section{Parameters, structure, and kernel types}
\label{sec:gen-typing-setup}

Fix the design $D=\{c_a\}_{a=1}^{M}$, the agent set $V=\{1,\dots,N\}$, and the measurable spaces
and sigma-finite reference measures of \Cref{ch:probability}.  Structural data
$X\in\mathsf X$ contain a finite directed acyclic graph $\Gamma=(V,E)$, design incidence, and the
geometric data. Writing $\operatorname{pa}(i)$ for the parents of $i$, the incidence is parent
closed: if $(i,a)\in\mathfrak I_D$ and $j\in\operatorname{pa}(i)$, then $(j,a)\in\mathfrak I_D$. Thus the induced graph
$\Gamma_a$ on $V_a$ contains every parent of each participating vertex and is acyclic. Numerical
parameters that a learning coordinate may move are denoted $\theta$. This split is declared rather
than derived; all typing requirements below apply to both symbols.
\status{DEFINITION}\label{def:gen-parameter-split}

Write $\mathcal R=\{r:\operatorname{pa}(r)=\varnothing\}$. At design point $c_a$, for each $i\in V_a$, declare normalized
Markov kernels
\begin{equation}
P_{\theta,i}^{m}(dm_i\given m_{\operatorname{pa}(i)},X),
\qquad
P_{\theta,i}^{k}(dk_i\given k_{\operatorname{pa}(i)},m_i,X),
\qquad
P_{\theta,i}^{o}(do_i\given k_i,m_i,X).
\label{eq:gen-kernels}
\end{equation}
At a root, the first factor is a proper model law and the second is a normalized bridge from the
same-agent model coordinate to the belief coordinate. The observation component is retained as a
normalized singleton-scope record mechanism; in the population construction below it is used only
through the owned record kernel $K_r$ to which it is assigned. \status{DEFINITION}

For each admissible incidence $(i,a)$, specialize the source and target
spaces of \Cref{def:prob-model-evaluation} by
\[
(\mathsf S_{i,a}^X,\mathscr S_{i,a}^X)
:=\left(
\prod_{j\in\operatorname{pa}(i)}\mathsf K_{j,a},
\bigotimes_{j\in\operatorname{pa}(i)}\mathscr K_{j,a}
\right),
\qquad
(\mathsf T_{i,a}^X,\mathscr T_{i,a}^X)
:=(\mathsf K_{i,a},\mathscr K_{i,a}),
\]
while retaining the already declared presentation space
\((\mathsf M_{i,a},\mathscr M_{i,a})\). The empty parent product is a
singleton with its trivial sigma-algebra. For fixed $(\theta,X)$, the
incidence evaluator used by the state factor is
\begin{equation}
\operatorname{ev}_{i,a}(m_i)=K^{X,a}_{i,m_i},
\qquad
K^{X,a}_{i,m_i}(k_{\operatorname{pa}(i)},B)
:=P_{\theta,i}^{k}(B\given k_{\operatorname{pa}(i)},m_i,X),
\label{eq:gen-evaluated-state-kernel}
\end{equation}
for every $m_i\in\mathsf M_{i,a}$, every complete parent-state tuple in the
displayed source, and $B\in\mathscr K_{i,a}$. Thus $m_i$ is the sampled
presentation supplied to the evaluator, while every parent-state coordinate
remains in the source of its normalized kernel. Together with the separate
model factor, which retains every parent-model coordinate, this specialization
reads the complete source block \(y_{\operatorname{pa}(i),a}\).
\status{DEFINITION}

For $i\in V_a$, write
$\mathsf Y_{\operatorname{pa}(i),a}=\prod_{j\in\operatorname{pa}(i)}\mathsf Y_{j,a}$.
The ordered composite of the model and state kernels is the normalized agent mechanism
\begin{equation}
\begin{split}
G_{\theta,i,a}&:
(\mathsf X\times\mathsf Y_{\operatorname{pa}(i),a})
\rightsquigarrow\mathsf Y_{i,a},\\
G_{\theta,i,a}\bigl(d(k_i,m_i)\given y_{\operatorname{pa}(i),a},X\bigr)
&:=P_{\theta,i}^{m}(dm_i\given m_{\operatorname{pa}(i)},X)
K^{X,a}_{i,m_i}(k_{\operatorname{pa}(i)},dk_i).
\end{split}
\label{eq:gen-agent-mechanism}
\end{equation}
Its conditioning coordinates precede its target block, and integration in $k_i$ followed by $m_i$
gives unit mass. The model factor retains the complete parent-model tuple and the evaluated state
factor retains the complete parent-state tuple, so their composition reads the full
$y_{\operatorname{pa}(i),a}$ source. Under the design-point conditional-independence hypothesis,
the finite-design
agent mechanism is
\begin{equation}
G_{\theta,i,D}:=\bigotimes_{a\in D_i}G_{\theta,i,a}.
\label{eq:gen-agent-design-mechanism}
\end{equation}
The two factors inside each $G_{\theta,i,a}$ are composed in their displayed dependency order.
The tensor product in $G_{\theta,i,D}$ is the product kernel across independent design points. \status{DEFINITION}

The channels are asymmetric by declaration. At a root incidence \((i,a)\), the
parent-state source in \eqref{eq:gen-evaluated-state-kernel} is the singleton
with element $\ast$. For
\(s\in\mathcal P(\mathsf M_{i,a},\mathscr M_{i,a})\), the evaluated kernel
induces the incidence-level map on laws
\begin{equation}
\widetilde\phi_{i,a}(s)(A)
=\int_{\mathsf M_{i,a}}K^{X,a}_{i,m}(\ast,A)s(dm)
=\int_{\mathsf M_{i,a}}P_{\theta,i}^{k}(A\given m,X)s(dm),
\qquad A\in\mathscr K_{i,a}.
\label{eq:gen-root-law-bridge}
\end{equation}
Only this pushforward-on-laws map can represent the design-point value of the
local associated-bundle morphism
$\widetilde\Phi:(\mathcal E_m)_{c_a}\to(\mathcal E_b)_{c_a}$, and only if its
contextual and gauge compatibility is declared. At a nonroot, the evaluated
kernel is indexed by the complete parent
belief tuple; it
does not define a map $\mathcal E_m\to\mathcal E_b$ unless that parent is fixed or the domain is
enlarged.  Otherwise it is a separately declared kernel-valued bridge.  In every case the
generative dependence runs from model samples to belief samples.  The geometric existence of
$\Phi:\mathcal E_b\to\mathcal E_m$ neither inserts a reverse factor into
\eqref{eq:gen-kernels} nor forces cyclic generative dependence.  A model using both directions
would need a separately normalized joint and cannot infer its factorization from the two
associated-bundle morphisms.

\section{Iterated composition and latent densities}
\label{sec:gen-composition}

For each $a$, let $N_a:=|V_a|$ and choose a topological ordering
$v_{a,1},\dots,v_{a,N_a}$ of $\Gamma_a$. The one-design-point latent law is the
successive kernel composition
\begin{equation}
\mathbb P_{\theta,V,a}^{Y}(dy_a\given X)
:=\overrightarrow{\prod_{\ell=1}^{N_a}}
G_{\theta,v_{a,\ell},a}
\bigl(dy_{v_{a,\ell},a}\given y_{\operatorname{pa}(v_{a,\ell}),a},X\bigr).
\label{eq:gen-composition}
\end{equation}
The arrow denotes ordered application of kernels, not multiplication of measures. For bounded
measurable $f$, its meaning is
\begin{equation}
\int f(y_a)\mathbb P_{\theta,V,a}^{Y}(dy_a\given X)
=\int G_{\theta,v_{a,1},a}(dy_{v_{a,1},a}\given X)\cdots
\int G_{\theta,v_{a,N_a},a}
\bigl(dy_{v_{a,N_a},a}\given y_{\operatorname{pa}(v_{a,N_a}),a},X\bigr)f(y_a).
\label{eq:gen-iterated-integral}
\end{equation}
Every conditioning coordinate has been introduced before its kernel is applied.

\medskip
\constructionheading{Finite directed latent law}{cons:gen-finite-directed-law}
Equation~\eqref{eq:gen-iterated-integral} defines a normalized measurable probability kernel
$X\mapsto\mathbb P_{\theta,V,a}^{Y}(\cdot\given X)$. \status{ESTABLISHED}
\medskip

\noindent\emph{Proof.}
Integrating a bounded measurable function against a probability kernel preserves boundedness and
measurability in the remaining variables. Finite induction along the topological order therefore
defines a measure and a measurable kernel in $X$; applying the induction to $f\equiv1$ gives total
mass one. This is the finite kernel construction underlying directed graphical models
\citep{Klenke2020,Lauritzen1996,Wainwright2008}. \hfill$\square$

Distinct design points are conditionally independent given $X$:
\begin{equation}
\mathbb P_{\theta,V}^{Y}(dy\given X)
=\bigotimes_{a=1}^{M}\mathbb P_{\theta,V,a}^{Y}(dy_a\given X).
\label{eq:gen-design-product}
\end{equation}
Equivalently, each agent contributes the bundled mechanism $G_{\theta,i,D}$ of
\eqref{eq:gen-agent-design-mechanism}. This modeling hypothesis excludes residual cross-design
dependence and is used in the finite-design normalization and factor count.
\status{HYPOTHESIS}\label{hyp:gen-design-product}

When the component kernels have densities relative to the declared component reference measures,
\begin{equation}
g_{\theta,i,a}(k_i,m_i\given y_{\operatorname{pa}(i),a},X)
=p_{\theta,i}^{m}(m_i\given m_{\operatorname{pa}(i)},X)
p_{\theta,i}^{k}(k_i\given k_{\operatorname{pa}(i)},m_i,X),
\qquad
p_{\theta,V,a}^{Y}(y_a\given X)=\prod_{i\in V_a}g_{\theta,i,a}.
\label{eq:gen-density}
\end{equation}
Consequently
\begin{equation}
p_{\theta,V}^{Y}(y\given X)=\prod_{a=1}^{M}p_{\theta,V,a}^{Y}(y_a\given X),
\qquad
\mathbb P_{\theta,V}^{Y}(dy\given X)
=p_{\theta,V}^{Y}(y\given X)\nu_D^Y(dy).
\label{eq:gen-design-density}
\end{equation}
The measure $\mathbb P_{\theta,V}^{Y}$ is authoritative; $p_{\theta,V}^{Y}$ denotes only its
Radon--Nikodym density.

\section{The typing prohibition}
\label{sec:gen-typing}

\requirementheading{Typing prohibition}{req:gen-typing-prohibition}
\emph{Each factor in \eqref{eq:gen-kernels} is a function only of $\theta$, $X$, and its displayed
parent and same-agent latent coordinates.  No generative kernel may take as an argument the
recognition law $\mathbb Q_{V,o,X}$, one of its local marginals or parameters, or the posterior $\boldsymbol\Pi_{\theta,V,o,X}$.}
\status{DEFINITION}

This requirement defines the fixed-joint construction and proves nothing.  A transition built from
a live peer recognition marginal is a map from recognition laws to kernels, not a factor of a
joint fixed before inference.

\propositionheading{No distinguished target}{prop:gen-no-distinguished-target}
Let $\mathcal Q$ be a nonempty family of recognition laws.  If a factor in
\eqref{eq:gen-composition} is replaced by a nonconstant $Q$-indexed family, producing a normalized
joint $P_{\theta,Q}$ for each $Q\in\mathcal Q$, assume at the displayed observation that
$0<p_{\theta,Q}(o\given X)<\infty$, that
$Q\ll P_{\theta,Q}(\cdot\given o,X)$, and that the log-density terms defining the ELBO and KL
divergence are integrable. Then for each fixed $Q$
\begin{equation}
\log p_{\theta,Q}(o\given X)
=\Lelbo(Q;X)
+\KL\left(Q\middle\Vert P_{\theta,Q}(\cdot\given o,X)\right).
\label{eq:gen-moving-target}
\end{equation}
The family itself selects no joint as ``the model.''  Writing
$e(Q)=\log p_{\theta,Q}(o\given X)$, two members satisfy
\begin{equation}
\begin{aligned}
\Lelbo(Q';X)-\Lelbo(Q;X)
={}&e(Q')-e(Q)\\
&-\left[
\KL\bigl(Q'\Vert P_{\theta,Q'}(\cdot\given o,X)\bigr)
-\KL\bigl(Q\Vert P_{\theta,Q}(\cdot\given o,X)\bigr)
\right].
\end{aligned}
\label{eq:gen-elbo-difference}
\end{equation}
Thus a divergence decrease certifies improvement relative to one fixed reference value exactly
when the two members lie on one level set of $e$. \status{ESTABLISHED}

\noindent\emph{Proof.}
Apply the exact ELBO identity to each fixed pair $(Q,P_{\theta,Q})$ and subtract the two
instances.  No member is distinguished because no selection rule was declared. \hfill$\square$

\propositionheading{The stronger consequences fail}{prop:gen-moving-target-witness}
Nonconstant dependence of $P_{\theta,Q}$ on $Q$ does not imply that $e$ varies, that no ascent can
improve a fixed quantity, or that no single member's evidence bounds the family of bounds.
\status{ESTABLISHED}

\noindent\emph{Proof.}
Take one binary latent $y$, one binary observation $o$, and
$Q_\beta=\operatorname{Ber}(\beta)$ for $\beta\in(0,1)$.  Let
$g(\beta)=\beta/2+1/4$ and define
\begin{equation}
P_{\theta,Q_\beta}(o,y\given X)
=\tfrac12\operatorname{Ber}\bigl(g(\beta)\bigr)(y).
\label{eq:gen-moving-target-witness}
\end{equation}
The joint is normalized and varies injectively with $\beta$, but
$e(Q_\beta)=\log(1/2)$ for every $\beta$.  The posterior is
$\operatorname{Ber}(g(\beta))$, so the gap at
$\beta=9/10,7/10,1/2$ is respectively
\begin{equation}
\tfrac9{10}\log\tfrac97+\tfrac1{10}\log\tfrac13
=0.1163218\ldots,
\qquad
\tfrac7{10}\log\tfrac76+\tfrac3{10}\log\tfrac34
=0.0216009\ldots,
\qquad
0.
\label{eq:gen-moving-target-values}
\end{equation}
The bound therefore increases to the fixed evidence, and the member at $\beta=1/2$ has evidence
equal to the supremum of all these bounds. \hfill$\square$

These are closed-form relative entropies, not numerical experiments.  The development declines the
three stronger statements as false in this witness, while retaining the absence of a distinguished
member. \status{NOT-CLAIMED}

The typing prohibition does not ban population coupling through latents inside
a fixed joint; it bans a generative factor from reading the posterior that the
joint is supposed to determine. \status{DEFINITION}

\section{Normalized record extension}
\label{sec:gen-normalization}

\propositionheading{Normalization by latent composition and record attachment}{prop:gen-exact-normalization}
Under the declared conditional independence of record coordinates given $(X,y)$, the full population law is the kernel
\begin{equation}
\mathbb P_{\theta,V}(dy,do\given X)
=
\mathbb P_{\theta,V}^{Y}(dy\given X)
\prod_{a\in\mathfrak A}K_a(X,y_{\partial a,D},do_a).
\label{eq:gen-record-extension}
\end{equation}
It is normalized for every admissible $(\theta,X)$:
\begin{equation}
\mathbb P_{\theta,V}(\mathsf Y_D\times\mathsf O_D\given X)=1.
\label{eq:gen-normalization}
\end{equation}
If the record kernels have jointly measurable densities
$K_a(X,y_{\partial a,D},do_a)=\ell_a(o_a\given X,y_{\partial a,D})\nu_a(do_a)$
and $\nu_D^O=\bigotimes_{a\in\mathfrak A}\nu_a$, then
\begin{equation}
p_{\theta,V}(o,y\given X)
=p_{\theta,V}^{Y}(y\given X)
\prod_{a\in\mathfrak A}\ell_a(o_a\given X,y_{\partial a,D})
\label{eq:gen-record-density}
\end{equation}
is its normalized density. \status{ESTABLISHED}

\noindent\emph{Proof.}
Measurability follows by finite kernel composition. For normalization, apply Tonelli's theorem and
integrate each record coordinate first. Every $K_a$ has unit mass in $o_a$, so their finite product
disappears and leaves $\mathbb P_{\theta,V}^{Y}(dy\given X)$. The latter has unit mass by the
finite topological-order construction and the finite design product. The same calculation with the
displayed Radon--Nikodym derivatives proves the density statement. \hfill$\square$

Cycles in the finite record hypergraph do not obstruct \eqref{eq:gen-record-extension}: record
coordinates are conditionally attached outputs of one already constructed latent law. By contrast,
reciprocal latent-state conditionals require a declared order or schedule, a common-cause
construction, or a globally normalized undirected model. Local normalization of those reciprocal
conditionals alone does not supply a joint.

\section{The normalized undirected alternative}
\label{sec:gen-undirected}

An undirected cyclic model is valid when its global normalizer exists.  For nonnegative node and
edge potentials, define
\begin{align}
\Psi_X(o,y)
&=\prod_{i\in V}\psi_i(o_i,k_i,m_i;X)
\prod_{\{i,j\}\in E}\psi_{ij}(k_i,m_i,k_j,m_j;X),
\label{eq:gen-gibbs-potential}\\
Z_X
&=\int\Psi_X(o,y)\nu_D^O(do)\nu_D^Y(dy),
\qquad 0<Z_X<\infty,
\label{eq:gen-gibbs-partition}\\
p_{\theta,V}^{\mathrm{Gibbs}}(o,y\given X)
&=Z_X^{-1}\Psi_X(o,y).
\label{eq:gen-gibbs-density}
\end{align}
Finiteness and positivity of $Z_X$ are model-existence conditions. \status{DEFINITION}

\propositionheading{Local normalization is insufficient}{prop:gen-gibbs-counterexample}
Normalized node potentials and an everywhere positive, finite, continuous edge potential need not
give finite $Z_X$. \status{ESTABLISHED}

\noindent\emph{Proof.}
On $\R^2$ take
\begin{equation}
\psi_i(y_i)=(2\pi)^{-1/2}e^{-y_i^2/2},
\qquad
\psi_{12}(y_1,y_2)=e^{cy_1y_2},
\qquad c\geq1.
\label{eq:gen-gibbs-counterexample}
\end{equation}
The quadratic precision is
\begin{equation}
\Lambda_{\mathrm{ex}}=\begin{pmatrix}1&-c\\-c&1\end{pmatrix},
\qquad
\operatorname{spec}\Lambda_{\mathrm{ex}}=\{1-c,1+c\}.
\label{eq:gen-gibbs-counterexample-matrix}
\end{equation}
In orthonormal sum-and-difference coordinates the integral factorizes, and the factor in the
$1-c$ direction diverges. \hfill$\square$

The document chooses the directed construction because its normalization does not depend on the
parameters, not because normalized cyclic Gaussian Markov random fields are invalid.

\section{Covariance under independent channel-frame re-expression}
\label{sec:gen-gauge}

At each admissible incidence $(i,a)\in\mathfrak I_D$, independently re-express the belief and
model frame coordinates by
$g_{a,i}^b=g_i^b(c_a)\in G$ and $g_{a,i}^m=g_i^m(c_a)\in G$.  Define the
old-to-new represented sample-coordinate maps
\begin{equation}
R_{a,i}^b:=\rho_b(g_{a,i}^b):\mathsf K\longrightarrow\mathsf K,
\qquad
R_{a,i}^m:=\rho_m(g_{a,i}^m):\mathsf M\longrightarrow\mathsf M.
\label{eq:gen-rep-actions}
\end{equation}
With the passive-section convention of
\eqref{eq:geo-local-reframing}, this sample-coordinate transformation
corresponds to the section rechoice
$u_i^{x\prime}=u_i^x\cdot(g_i^x)^{-1}$, whose value at $c_a$ uses
$(g_{a,i}^x)^{-1}$.  Thus $R_{a,i}^x$ acts within the
channel's sample fiber, while the group element $(g_i^x)^{-1}\in G$, not the
represented operator $(R_{a,i}^x)^{-1}$, multiplies the principal frame on the right.  This
distinction remains necessary when the representation is not faithful.
Within a formula for one design point, suppress the index $a$ on the participating vertices and write $g_i^x,R_i^x$.
They act on latent coordinates as $k_i'=R_i^bk_i$ and $m_i'=R_i^mm_i$, while observations are
gauge inert. \status{DEFINITION}

To permit all such independent coordinate changes, set
\(D_{ij}:=D_i\cap D_j\) for each edge \((i,j)\). The finite-design structural
data carry graph-link copies \(\Theta_{a,ij}^x\) exactly for \(a\in D_{ij}\).
At fixed \(a\in D_{ij}\) write \(\Theta_{ij}^x\); their discrete vertex frame
changes are
\begin{equation}
h_{a,i}^x:=g_{a,i}^x=g_i^x(c_a),
\qquad
a_{a,i}^x:=(h_{a,i}^x)^{-1},
\label{eq:gen-contextual-frame-change}
\end{equation}
in the notation of \eqref{eq:geo-regime-two-gauge-law}.  The contextual coordinate element
$h_{a,i}^x$, abbreviated $h_i^x$ at one design point, is unrelated to the relative
principal-frame field $h_i$ defined by \eqref{eq:geo-relative-frame}.  The suppressed
one-design-point law is
\begin{equation}
\Theta_{ij}^{b\prime}=h_i^b\Theta_{ij}^b(h_j^b)^{-1},
\qquad
\Theta_{ij}^{m\prime}=h_i^m\Theta_{ij}^m(h_j^m)^{-1}.
\label{eq:gen-gauge-links}
\end{equation}
Consequently the represented links obey
\begin{equation}
\rho_x(\Theta_{ij}^{x\prime})
=R_i^x\rho_x(\Theta_{ij}^x)(R_j^x)^{-1},
\qquad x\in\{b,m\},
\label{eq:gen-represented-gauge-links}
\end{equation}
and the inverse-orientation rule is preserved.  Equations
\eqref{eq:geo-regime-two-gauge-law} and \eqref{eq:gen-gauge-links} are therefore the same passive
law written with mutually inverse rechoice conventions. \status{ESTABLISHED}

If instead one structural link is constrained to be identical at every point of
\(D_{ij}\), an admissible frame rechoice must preserve that constraint:
for every \(a,a'\in D_{ij}\) and every constrained edge \((i,j)\),
\begin{equation}
h_{a,i}^x\Theta_{ij}^x(h_{a,j}^x)^{-1}
=h_{a',i}^x\Theta_{ij}^x(h_{a',j}^x)^{-1}.
\label{eq:gen-shared-link-admissibility}
\end{equation}
\status{DEFINITION}

At one fixed shared-link datum, define the admissible family to be the solution set of
\eqref{eq:gen-shared-link-admissibility}.  Define the group that acts on the entire constrained
submodel to be the setwise stabilizer of the shared-link constraint set.  A larger
transformation law requires an explicitly declared parameter compensation.
\status{DEFINITION}

\propositionheading{Shared-link admissibility need not be closed under composition}{prop:gen-shared-link-nonclosure}
For $G=\GL(2,\R)$, two design points at which both endpoints participate, one constrained edge,
and $\Theta_{ij}=I$, the solution set of \eqref{eq:gen-shared-link-admissibility} is not closed under
pointwise multiplication. \status{ESTABLISHED}

\noindent\emph{Proof.}
Let $S=\operatorname{diag}(2,1)$ and
$C=\begin{psmallmatrix}1&1\\0&1\end{psmallmatrix}$.  The rechoice
$h_{1,i}=h_{1,j}=I$, $h_{2,i}=h_{2,j}=S$ sends the shared link to $I$ at both design points.
The rechoice $\widetilde h_{a,i}=C$, $\widetilde h_{a,j}=I$ sends it to $C$ at both points.
For their pointwise product $h\widetilde h$, the two transformed links are $C$ and
$SCS^{-1}=\begin{psmallmatrix}1&2\\0&1\end{psmallmatrix}$, respectively, so the product violates
\eqref{eq:gen-shared-link-admissibility}. \hfill$\square$

Context-independent rechoices form a sufficient subgroup, but the first
rechoice in the proof shows that they are not necessary. \status{ESTABLISHED}

At the unconstrained finite-design level, declaring the coordinate re-expressions independently
gives the passive product $\prod_{(i,a)\in\mathfrak I_D}(G\times G)$. If they must arise as values of smooth maps
$g_i^x:\mathcal C_i\to G$, the realized group is the image of the corresponding
incidence-restricted evaluation map; it equals the displayed product whenever that evaluation map
is surjective. For example, this holds when each relevant $\mathcal C_i$ is smooth and
positive-dimensional, the points $\{c_a:a\in D_i\}$ are distinct, and $G$ is connected, by smooth
interpolation in disjoint coordinate neighborhoods. This product records two independently chosen
trivializations of associated bundles induced from the same $\mathscr P_G$; it is not a second principal
structure group. The active principal gauge group is $\operatorname{Aut}_G(\mathscr P_G)$ over
$\operatorname{id}_{\mathcal C}$; one automorphism acts on both associated bundles, with local
functions related by \eqref{eq:geo-diagonal-gauge-functions}.  Its image in the passive product is
contained in the pointwise $h_i$-twisted diagonal and is contained in the ordinary diagonal in the
common-frame specialization $h_i=e_G$.  Equality with the corresponding pointwise diagonal holds
only when the relevant active-gauge evaluation map is surjective; disconnected structure groups
or global bundle constraints can obstruct that surjectivity.  It is not the principal right
$G$-action itself. A
context-independent frame re-expression is the subgroup for which $g_{a,i}^x$ is constant as
$a$ ranges over $D_i$.
\status{DEFINITION}\label{def:gen-product-action}

\medskip
\hypothesisheading{Kernel covariance}{hyp:gen-kernel-covariance}
The transformed parameters $(\theta',X')$ are declared so that every nonroot factor obeys
\begin{align}
P_{\theta',i}^{m}
\left(\cdot\given R_{\operatorname{pa}(i)}^mm_{\operatorname{pa}(i)},X'\right)
&=(R_i^m)_{\#}
P_{\theta,i}^{m}
\left(\cdot\given m_{\operatorname{pa}(i)},X\right),
\label{eq:gen-gauge-pushforward-model}\\
P_{\theta',i}^{k}
\left(\cdot\given R_{\operatorname{pa}(i)}^bk_{\operatorname{pa}(i)},R_i^mm_i,X'\right)
&=(R_i^b)_{\#}
P_{\theta,i}^{k}
\left(\cdot\given k_{\operatorname{pa}(i)},m_i,X\right),
\label{eq:gen-gauge-pushforward-state}\\
P_{\theta',i}^{o}
\left(\cdot\given R_i^bk_i,R_i^mm_i,X'\right)
&=P_{\theta,i}^{o}
\left(\cdot\given k_i,m_i,X\right),
\label{eq:gen-gauge-pushforward-obs}
\end{align}
with the analogous root identities. Every interaction-record kernel also obeys
\begin{equation}
K_a\left(X',R_{\partial a,D}y_{\partial a,D},\cdot\right)
=K_a\left(X,y_{\partial a,D},\cdot\right),
\qquad a\in\mathfrak A,
\label{eq:gen-gauge-record-kernel}
\end{equation}
where $R_{\partial a,D}$ is the product restriction of $R$ to the incident agent blocks. Record coordinates are gauge inert. \status{HYPOTHESIS}
\medskip

The second identity is the covariance requirement for the model-to-belief sample
kernel. At each root incidence it implies that the induced law map
\eqref{eq:gen-root-law-bridge} obeys the design-point form of
\eqref{eq:geo-tildephi-gauge-law}; at a nonroot it remains a parent-indexed
kernel law unless an enlarged cross-bundle domain is declared. Neither route
uses $\Phi$.

\propositionheading{Joint measure covariance and evidence invariance}{prop:gen-product-evidence-invariance}
Let
$T(y,o)=(Ry,o)$, where
$R=\bigoplus_{(i,a)\in\mathfrak I_D}(R_{a,i}^b\oplus R_{a,i}^m)$. Under
\Cref{hyp:gen-kernel-covariance},
\begin{equation}
\mathbb P_{\theta',V}(\cdot\given X')=T_{\#}\mathbb P_{\theta,V}(\cdot\given X),
\label{eq:gen-gauge-joint-pushforward}
\end{equation}
as probability measures, without a density or reference-measure hypothesis.  Because $T$ is the
identity on observations, the observation marginals agree.  Relative to the common observation
reference measure this gives
\begin{equation}
p_{\theta',V}(o\given X')=p_{\theta,V}(o\given X)
\quad\text{for }\nu_D^O\text{-almost every }o.
\label{eq:gen-gauge-evidence-invariance}
\end{equation}
If latent densities exist and $R_\#\nu_D^Y\ll\nu_D^Y$, define
\begin{equation}
j_R(y'):=\frac{d(R_\#\nu_D^Y)}{d\nu_D^Y}(y').
\label{eq:gen-gauge-reference-rn}
\end{equation}
Then the density statement is
\begin{equation}
p_{\theta',V}(o,y'\given X')
=p_{\theta,V}(o,R^{-1}y'\given X)j_R(y')
\quad\text{for }(\nu_D^Y\otimes\nu_D^O)\text{-almost every }(y',o).
\label{eq:gen-gauge-rn-density}
\end{equation}
On the further differentiable equal-dimensional Euclidean tier, when every represented block is
an invertible linear map and $\nu_D^Y$ is product Lebesgue measure,
$j_R=|\det R|^{-1}$ and hence
\begin{equation}
p_{\theta',V}(o,y'\given X')
=|\det R|^{-1}p_{\theta,V}(o,R^{-1}y'\given X).
\label{eq:gen-gauge-lebesgue-density}
\end{equation}
Equal dimensionality here is between the source and target coordinates within each channel; it
does not impose $K=d_m$.  The determinant formula is not asserted for mixed or atomic reference
measures.
\status{ESTABLISHED}

\noindent\emph{Proof.}
Apply \eqref{eq:gen-gauge-pushforward-model} and
\eqref{eq:gen-gauge-pushforward-state} along the topological ordering of each $\Gamma_a$.
Finite kernel induction gives
\[
\mathbb P_{\theta',V,a}^{Y}(\cdot\given X')
=R_{a\#}\mathbb P_{\theta,V,a}^{Y}(\cdot\given X),
\]
where $R_a$ is the product restriction of $R$ to the block at design point $c_a$,
and the product over design points in \eqref{eq:gen-design-product} therefore gives
$\mathbb P_{\theta',V}^{Y}(\cdot\given X')=R_\#\mathbb P_{\theta,V}^{Y}(\cdot\given X)$.
Now let $F$ be a bounded measurable function on the joint latent--observation space. Applying
\eqref{eq:gen-record-extension}, then the latent pushforward, gives
\[
\begin{split}
\int F d\mathbb P_{\theta',V}
&=\int\mathbb P_{\theta,V}^{Y}(dy\given X)
\int F(Ry,o)\prod_{r\in\mathfrak A}
K_r\left(X',R_{\partial r,D}y_{\partial r,D},do_r\right)\\
&=\int\mathbb P_{\theta,V}^{Y}(dy\given X)
\int F(Ry,o)\prod_{r\in\mathfrak A}
K_r\left(X,y_{\partial r,D},do_r\right)\\
&=\int F\circ T d\mathbb P_{\theta,V},
\end{split}
\]
where the second equality applies \eqref{eq:gen-gauge-record-kernel} to every record factor. This
proves \eqref{eq:gen-gauge-joint-pushforward} without using the legacy singleton observation
kernel $P_{\theta,i}^{o}$. Projection to the observation coordinate now shows that the observation
measures agree, and their densities satisfy \eqref{eq:gen-gauge-evidence-invariance}. The joint
measure identity and the Radon--Nikodym representations give
\[
\int F(y',o)p_{\theta',V}(o,y'\given X')\nu_D^Y(dy')\nu_D^O(do)
=\int F(Ry,o)p_{\theta,V}(o,y\given X)\nu_D^Y(dy)\nu_D^O(do).
\]
Rewriting the right side against $R_\#\nu_D^Y$ and applying
\eqref{eq:gen-gauge-reference-rn} gives \eqref{eq:gen-gauge-rn-density}. On the stated Lebesgue
tier, linear change of variables gives
$d(R_\#\nu_D^Y)/d\nu_D^Y=|\det R|^{-1}$, proving
\eqref{eq:gen-gauge-lebesgue-density}. \hfill$\square$

Evidence invariance is a redundancy constraint, not an empirical prediction: the action changes
latent coordinates while leaving the observation law fixed.  The linear-Gaussian chapter verifies
\Cref{hyp:gen-kernel-covariance} parameter by parameter and then displays the stronger
common-frame reduction.

The passive coordinate action on the full parameter space can have nontrivial
stabilizers, so its orbit space need not be a manifold. Subtracting the
dimension of the coordinate group from a parameter count is unjustified
without freeness and properness. Determining generic stabilizers and a regular
quotient remains open; nothing later depends on that quotient.
\status{OPEN}\label{open:gen-gauge-quotient}
